% !PRINTED 230
\weberpage{230}
%<<CONT
\weq{8}{\begin{aligned}
(-1)^{\frac{n^2-1}{8}}\frac{2^{\frac{n-1}{3}}P_{00}^2P_{11}}{\sqrt{\varkappa\varkappa'}^{\frac{n-1}{3}}} &= \mathop{\Pi}\limits_{1,\frac{n-1}{2}}^{h}\frac{\operatorname{sn}h\Omega\,\operatorname{dn}h\Omega}{\operatorname{cn}h\Omega},\\
(-1)^{\frac{n^2-1}{8}}\frac{2^{\frac{n-1}{3}}P_{10}^2P_{11}}{\sqrt{\varkappa\varkappa'}^{\frac{n-1}{3}}} &= \mathop{\Pi}\limits_{1,\frac{n-1}{2}}^{h}\frac{\operatorname{sn}h\Omega\,\operatorname{cn}h\Omega}{\operatorname{dn}k\Omega},\\ % [sic] printed \textquotedbl{}dn k Omega\textquotedbl{} for \textquotedbl{}dn h Omega\textquotedbl{}
(-1)^{\frac{n^2-1}{8}}\frac{2^{\frac{n-1}{3}}P_{01}^2P_{11}}{\sqrt{\varkappa\varkappa'}^{\frac{n-1}{3}}} &= \mathop{\Pi}\limits_{1,\frac{n-1}{2}}^{h}\frac{\operatorname{sn}h\Omega}{\operatorname{cn}h\Omega\,\operatorname{dn}h\Omega}.
\end{aligned}}

Ein wichtiges Resultat ergibt sich aber noch durch Anwendung des Multiplikationstheorems. Wenn man in der letzten Formel II des § 57 $u = 2h\varpi$ setzt, so findet man
\weq{9}{e^{4\pi i\omega h^2\mu'^2}\vartheta_{01}(2h\varpi)^{n^2} = \frac{\vartheta_{01}^{n^2}}{D(\operatorname{sn}^2h\Omega)},}
worin, wie wir uns erinnern, $D$ eine ganze rationale Funktion ist, deren Koeffizienten rational aus $\varkappa^2$ und ganzen Zahlen gebildet sind.

Nehmen wir das Produkt aus diesen Ausdrücken, für $h = 1, 2, \ldots, \frac{n-1}{2}$, so folgt aus der letzten Gleichung (2) (weil bekanntlich $\Sigma h^2 = \frac{n(n^2-1)}{24}$):
\weq{10}{P_{01}^{n^2} = \mathop{\Pi}\limits_{1,\frac{n-1}{2}}^{h}\frac{1}{D(\operatorname{sn}^2h\Omega)}.}

Wir schließen hieraus, daß auch $P_{01}^{n^2}$ die Wurzel einer Transformationsgleichung ist. Dies ist nichts Neues, wenn $n$ durch 3 teilbar ist, wohl aber, wenn $n$ nicht durch 3, also $n^2 - 1$ durch 3 teilbar ist. Wir erhalten dann nämlich aus (10) mit Benutzung von (5), (7), (8):
\weq{11}{\begin{aligned}
P_{01} &= 2^{\frac{(n-1)(n^2-1)}{6}}\mathop{\Pi}\limits_{1,\frac{n-1}{2}}^{h}\frac{(\operatorname{cn}h\Omega)^{\frac{n^2-1}{3}}(\operatorname{dn}h\Omega)^{\frac{n^2-1}{3}}}{D(\operatorname{sn}^2h\Omega)},\\
P_{00} &= 2^{\frac{(n-1)(n^2-1)}{6}}\mathop{\Pi}\limits_{1,\frac{n-1}{2}}^{h}\frac{(\operatorname{cn}h\Omega)^{\frac{n^2-1}{3}}(\operatorname{dn}h\Omega)^{\frac{n^2+2}{3}}}{D(\operatorname{sn}^2h\Omega)},\\
P_{10} &= 2^{\frac{(n-1)(n^2-1)}{6}}\mathop{\Pi}\limits_{1,\frac{n-1}{2}}^{h}\frac{(\operatorname{cn}h\Omega)^{\frac{n^2+2}{3}}(\operatorname{dn}h\Omega)^{\frac{n^2-1}{3}}}{D(\operatorname{sn}^2h\Omega)},
\end{aligned}}
%CONT>>
